\documentclass[10pt]{article}

\usepackage[margin=1in]{geometry}
\usepackage{microtype}
\usepackage{xcolor}
\definecolor{linkblue}{RGB}{0, 70, 140}
\usepackage[colorlinks=true, urlcolor=linkblue, linkcolor=linkblue, citecolor=linkblue]{hyperref}

\setlength{\parskip}{6pt}
\setlength{\parindent}{0pt}
\pagestyle{empty}

\begin{document}

\begin{center}
 {\large\bfseries Shape-Constrained Boosted Symbolic Regression\\
    for Interpretable Neural Scaling Laws: A Cross-Testbed Investigation}\\[6pt]
    {\normalsize Project Lead: \href{https://krichelj.github.io/}{Joshua Shay Kricheli}\\
    PhD Candidate in Electrical Engineering and Computer Science, Syracuse University}
\end{center}

\smallskip
\noindent\textbf{Abstract / Summary.}
Training a modern AI language model can cost millions of dollars, so researchers use scaling laws to predict performance from model size and training data before committing resources. These equations guide major hardware and budget decisions, but their mathematical form is usually chosen manually and then fitted to measured results.

Our project uses symbolic regression to discover the scaling-law equation directly from data. Symbolic regression searches many candidate formulas and returns an interpretable equation rather than a black-box model. Unconstrained searches, however, often produce formulas that fit existing measurements but extrapolate unrealistically, predicting negative error or worse performance with more data.

We therefore require every candidate formula to satisfy admissibility axioms: larger models and additional data cannot increase error; improvements must diminish with scale; error must remain positive; and performance must approach an irreducible floor. We mathematically certify these properties across the full domain of interest, not only at measured points.

The formula is built incrementally from a standard scaling law, adding one correction at a time while preserving all guarantees. We also prove that the corrections improve fit without changing the irreducible error floor or the asymptotic rate of improvement. One search method rejects inadmissible corrections; another penalizes violations and drives them toward zero while bounding any remaining shortfall.

The project has two computational bottlenecks. Model training will use Cerebras CS-3 systems on PSC's Neocortex platform and will be compared with completed NVIDIA GPU runs. Formula discovery is a large CPU-bound search dominated by interval-arithmetic certification. Because it is highly parallel and relies mainly on scalar or imperfectly vectorized Python and JAX code, Stony Brook's high-core-count AMA27 AmpereOne cluster is well suited to the workload.

Deliverables will include a certified scaling-law formula fitted to our training data, an open-source discovery library, and a concise methods report. An initial set of trained models is already public.

\medskip
\noindent\textbf{Target System:} Neocortex and AMA27\\
\textbf{Field of science:} Machine Learning and Artificial Intelligence, specifically the foundations of large language models (neural scaling laws) and interpretable machine learning via symbolic regression.\\
\textbf{Key Requirements:} Cerebras wafer-scale systems and Arm CPUs

\vspace{8pt}
{\footnotesize Condensed abstract circulated to ByteBoost~2.0 participants (Dave Carlson, Stony Brook).}

\end{document}
